\begin{abstract}
The peptides displayed by \plainMHCI{} molecules can differ among tissues even when the antigen-presenting allele (MHC allele) and source protein are the same. These differences may reflect protein availability, antigen processing, cell composition, and experimental sampling, but most existing predictors focus on general binding to or presentation by an \plainMHC{} rather than tissue-associated preference. We therefore formulate a biological ranking problem (relative prediction task): given two peptides from the same parent protein under the same antigen-presenting restriction (MHC restriction), determine which has stronger presentation evidence in a specified tissue. The human benchmark contains 157 tissue--human-leukocyte-antigen combinations (tissue--HLA tasks) and the mouse benchmark contains 24 tissue--mouse-antigen-presenting-system combinations (tissue--H2 tasks), with both requiring more than 200 peptide pairs sharing the restriction and parent protein in each combination (matched pairs per task). TissuePMHC combines representations learned globally across human tissue--\plainHLA{} combinations (human tasks) with representations restricted to one \plainHLA{} allele. On the \plainFixedTest{}s, it reaches equal-weight across-task AUROC/AUPRC values of 0.8448/0.8348 in human, while a separately trained mouse model reaches 0.8562/0.8506. Under \plainPeptideOOF{}, AUROC decreases to 0.7652 in human and 0.7529 in mouse. Thus, both species retain predictive signal for previously represented tissue--antigen-presenting-molecule combinations (tissue--MHC combinations), but peptide reuse leads to a more optimistic estimate from familiar tasks (standard estimate). Removing tissue identity causes a further substantial decrease, showing that general peptide--antigen-presenting-molecule compatibility (peptide--MHC compatibility) alone does not reproduce the full result. However, the complete human and mouse architectures do not consistently outperform their strongest simpler controls after peptide separation. The benchmark therefore supports tissue-conditioned peptide ranking and quantifies the performance decrease for peptide identities absent from all training tasks (unseen-peptide generalization gap); it does not establish a causal tissue mechanism, performance for new tissues or antigen-presenting alleles (MHC alleles), independent-cohort validity, or clinical utility.
\end{abstract}
